\subsection{Sistemas de primer orden}

Un sistema de primer orden con retardo se describe mediante la función de
transferencia en la \cref{eq:transfer-function-FO}, donde \(k_p\) es la
ganancia, \(T_p\) la constante de tiempo y \(T_d\) el retardo.
\begin{equation}\label{eq:transfer-function-FO}
	G(s) = \frac{k_p}{T_p s + 1} e^{-T_d s}
\end{equation}

Para una entrada escalón unitario de amplitud \(A\) activada en \(b\),
la salida en el dominio de la frecuencia está dada por la
\cref{eq:output-s}. La transformada inversa de Laplace conduce a la
salida temporal, expresada en la \cref{eq:output-t}. Si la salida tiene
un desplazamiento (offset), este puede compensarse redefiniendo
\(y(t) = \tilde{y}(t) - \text{offset}\).
\begin{gather}
	Y(s) = \frac{A k_p}{s(T_p s + 1)} e^{-(T_d + b)s} \label{eq:output-s} \\
	y(t) = A k_p u(t - (T_d + b))
	\left[1 - e^{-(t - (T_d + b))T^{-1}_{p}} \right] \label{eq:output-t}
\end{gather}

Aquí, \(u(t)\) representa la función de Heaviside.

\subsection{Sistemas de segundo orden}

Para un sistema de segundo orden con retardo, la función de transferencia
incorpora dos constantes de tiempo, \(T_{p_1}\) y \(T_{p_2}\), como se
muestra en la \cref{eq:transfer-function-SO}.
\begin{equation}\label{eq:transfer-function-SO}
	G(s) = \frac{k_p}{(T_{p_1} s + 1) (T_{p_2} s + 1)} e^{-T_d s}
\end{equation}

Bajo la misma entrada escalón, las expresiones de la salida en los dominios
de la frecuencia y el tiempo se presentan en las
\cref{eq:output-s2,eq:output-t2}.
\begin{align}
	Y(s) = & \, A k_p \Bigg[\frac{1}{s}
	- \frac{T_{p_1}^2}{(T_{p_1} - T_{p_2})(T_{p_1} s + 1)} \nonumber \\
	& \, + \frac{T_{p_2}^2}{(T_{p_1} - T_{p_2})(T_{p_2} s + 1)}
	\Bigg] e^{-(T_d + b)s}
	\label{eq:output-s2}
\end{align}

\begin{align}
	y(t) = & \, A k_p u(t - (T_d + b)) \Bigg[1
	- \frac{T_{p_1}}{T_{p_1} - T_{p_2}} \, e^{-(t - (T_d + b))T^{-1}_{p_1}} \nonumber \\
	& \, + \frac{T_{p_2}}{T_{p_1} - T_{p_2}} \, e^{-(t - (T_d + b))T^{-1}_{p_2}}
	\Bigg]
	\label{eq:output-t2}
\end{align}
